% GENERATED FILE. Edit canonical lessons, then run npm run generate:books.
% canonical-source-hash: fnv1a64:c14de0e50c52fb0d
% canonical-lessons: ES-C57-ll

\chapter{Cl-, Pl-, Fl- to Ll-}
\label{ch:sound-cluster-ll}
\hlchaptermodality{29}

\begin{chapteropening}
\textbf{By the end of this chapter:} \emph{I can try three Latin clusters on an unfamiliar ll- word and find the English cousin.}

Complete ES-C57-ll: name llamar's cousins, give the three clusters, and decode llave.
\end{chapteropening}

\section[cl-, pl-, fl- to ll-]{cl-, pl-, fl- $\to$ ll-}

\label{lesson:ES-C57-ll}



\begin{quote}
Say \textquotedblleft{}my name is.\textquotedblright{} (\emph{Me llamo}.) That \emph{ll-} is the second sound law of this book, and you have been using it since your third chapter without being told.
\end{quote}

\begin{grammarlens}[title={Grammar Lens: three clusters, one destination}]
\begin{quote}
\textbf{Latin \emph{cl-}, \emph{pl-} and \emph{fl-} all became Spanish \emph{ll-}.}
\end{quote}

\emph{Llamar} is Latin \textbf{clāmāre}, \textquotedblleft{}to cry out.\textquotedblright{} English took the same root untouched: \textbf{claim}, \textbf{clamour}, \textbf{proclaim}, \textbf{exclaim}.

\begin{quote}
\emph{llamar} $\leftarrow$ \emph{clāmāre} $\to$ \textbf{claim}, \textbf{clamour}
\end{quote}

And it happened to all three clusters, which is what makes it a law rather than a one-off:

\noindent
\begin{tabularx}{\linewidth}{@{}>{\raggedright\arraybackslash}X>{\raggedright\arraybackslash}X>{\raggedright\arraybackslash}X@{}}
\toprule
\textbf{Spanish} & \textbf{from} & \textbf{English cousin} \\
\midrule
\emph{llamar} & \emph{clāmāre} & \textbf{claim}, \textbf{clamour} \\
\emph{llave} & \emph{clāvem} & \textbf{clavicle}, \textbf{clef}, \textbf{conclave} \\
\emph{llover} & \emph{plovere} & \textbf{pluvial} \\
\emph{lleno} & \emph{plēnum} & \textbf{plenty}, \textbf{plenary} \\
\emph{llano} & \emph{plānum} & \textbf{plain}, \textbf{plane} \\
\emph{llama} & \emph{flamma} & \textbf{flame}, \textbf{inflammable} \\
\bottomrule
\end{tabularx}

Three different starting points arriving at one sound is unusual, and it is useful precisely because it is untidy: an unfamiliar Spanish word in \emph{ll-} has \textbf{three} English shapes to try, and one of them is very often right.
\end{grammarlens}

\begin{culture}
You now hold three sound laws, and they are not three facts. They are three instances of the same thing.

\begin{itemize}
\raggedright
  \item \emph{-ct-} $\to$ \emph{-ch-} — \emph{noctem} $\to$ \emph{noche}, and English kept \textbf{nocturnal}.
  \item \emph{f-} $\to$ \emph{h-} — \emph{fabulārī} $\to$ \emph{hablar}, and English kept \textbf{fable}.
  \item \emph{cl-}, \emph{pl-}, \emph{fl-} $\to$ \emph{ll-} — \emph{clāmāre} $\to$ \emph{llamar}, and English kept \textbf{clamour}.
\end{itemize}

Every one has the same shape: \textbf{Spanish inherited the word and wore it down; English borrowed it later and left it closer to the Latin.} That is why the English cousin so often looks more Latin than the Spanish does, which feels backwards until you see why.

So the laws are not trivia about the past. They are a \textbf{conversion table} between the Spanish in front of you and the English already in your head, and they will keep working on words no course will ever teach you.
\end{culture}

\subsection*{Your turn}
\begin{itemize}
\raggedright
  \item \emph{Say it:} \textquotedblleft{}llamar\textquotedblright{} then \textquotedblleft{}clamour\textquotedblright{}; \textquotedblleft{}llave\textquotedblright{} then \textquotedblleft{}clavicle\textquotedblright{}; \textquotedblleft{}llama\textquotedblright{} then \textquotedblleft{}flame\textquotedblright{}
\end{itemize}

\begin{tcolorbox}[breakable,colback=teal!4,colframe=teal!35!black,title={Before you move on}]
Which three clusters became \emph{ll-}? (\emph{\textbf{cl-}}, \emph{\textbf{pl-}}, \emph{\textbf{fl-}}.) What is \emph{llamar}'s English cousin? (\textbf{Claim}.) And why does the English usually look more Latin? (Spanish \textbf{inherited and wore it down}; English \textbf{borrowed it later}.)
\end{tcolorbox}
